\chapter{Continuous Random Variables}
\label{ch:cont}

A waiting time, a height, a measurement error: these take values in an interval
of the real line, not in a countable set. Probability must now be spread over a
continuum, so that individual points get mass zero and intervals get positive
mass. The tool is the \emph{density}, and sums are replaced by integrals. This
chapter develops the theory and the standard continuous distributions, with
proofs.

\section{Densities}

\begin{definition}[Continuous random variable and density]
A random variable $X$ is \emph{continuous} (more precisely, \emph{absolutely
continuous}) if there is a function $f=f_X:\R\to[0,\infty)$, called its
\emph{probability density function} (PDF), such that
\[
F_X(x)=\Prob(X\le x)=\int_{-\infty}^xf(t)\,dt\qquad\text{for all }x\in\R.
\]
Necessarily $\int_{-\infty}^\infty f(t)\,dt=1$.
\end{definition}

\begin{proposition}\label{prop:density}
Let $X$ have density $f$. Then
\begin{enumerate}[label=(\roman*)]
\item $\Prob(a<X\le b)=\int_a^bf(t)\,dt$, and $\Prob(X\in B)=\int_Bf(t)\,dt$ for
every ``reasonable'' set $B$ (finite unions of intervals, and their limits);
\item $\Prob(X=x)=0$ for every $x$; hence $\Prob(a\le X\le b)=\Prob(a<X<b)$;
\item $F$ is continuous, and $F'(x)=f(x)$ at every point $x$ where $f$ is
continuous.
\end{enumerate}
\end{proposition}
\begin{proof}
(i) $\Prob(a<X\le b)=F(b)-F(a)$ by Theorem~\ref{thm:cdf}(iv). (ii) By
Theorem~\ref{thm:cdf}(v) the mass at $x$ equals the jump of $F$, which is zero
since an integral over an interval of shrinking length $\to0$ tends to $0$ (the
integrand is integrable). (iii) The fundamental theorem of calculus.
\end{proof}

\begin{nuance}{a density is not a probability}
The value $f(x)$ is \emph{not} $\Prob(X=x)$ (that is $0$), and $f$ can exceed $1$:
the uniform density on $[0,\tfrac12]$ equals $2$. The correct interpretation is
\[
\Prob(x<X\le x+h)=\int_x^{x+h}f(t)\,dt\approx f(x)\,h\quad(h\to0),
\]
so $f(x)$ is ``probability per unit length'' near $x$. Changing $f$ at a single
point (or on a set of length zero) does not change any probability, so a density
is determined only up to such changes.
\end{nuance}

\begin{nuance}{not every distribution without atoms has a density}
``Continuous CDF'' and ``has a density'' are different. The \emph{Cantor
distribution} (the law of $\sum_{k\ge1}2\epsilon_k3^{-k}$ with $\epsilon_k$ fair
Bernoulli) has a continuous CDF, no atoms, but is concentrated on the Cantor set
(length zero), so no density exists. Moreover, a distribution can be \emph{mixed},
having both atoms and a density part: for example $\Prob(X=0)=\tfrac12$ and
otherwise exponential. The CDF handles all cases uniformly, and
Chapter~\ref{ch:measure} gives the complete classification (Lebesgue
decomposition). In this chapter ``continuous'' always means ``has a density''.
\end{nuance}

\section{Expectation}

\begin{definition}[Expectation]
If $X$ has density $f$ and $\int\abs xf(x)\,dx<\infty$, then
$\Exp[X]=\int_{-\infty}^\infty xf(x)\,dx$. Variance, covariance and moments are
defined exactly as before: $\Var(X)=\Exp[(X-\mu)^2]$.
\end{definition}

\begin{theorem}[Tail-integral formula]\label{thm:tailint}
If $X\ge0$ has a density, then $\Exp[X]=\int_0^\infty\Prob(X>t)\,dt$.
\end{theorem}
\begin{proof}
Write $x=\int_0^x1\,dt=\int_0^\infty\indic_{\{t<x\}}\,dt$. Then
\[
\Exp[X]=\int_0^\infty x f(x)\,dx=\int_0^\infty\!\!\int_0^\infty\indic_{\{t<x\}}f(x)\,dt\,dx
=\int_0^\infty\!\!\int_t^\infty f(x)\,dx\,dt=\int_0^\infty\Prob(X>t)\,dt.
\]
The interchange of the order of integration is legitimate because the integrand
is non-negative (Tonelli's theorem, Chapter~\ref{ch:measure}).
\end{proof}

\begin{theorem}[LOTUS, continuous case]\label{thm:lotuscont}
If $X$ has density $f$ and $g\ge0$ (or $\int\abs{g}f<\infty$), then
$\Exp[g(X)]=\int g(x)f(x)\,dx$.
\end{theorem}
\begin{proof}
For $g\ge0$, apply Theorem~\ref{thm:tailint} to the non-negative variable
$g(X)$ and interchange integrals again:
\[
\Exp[g(X)]=\int_0^\infty\Prob(g(X)>t)\,dt=\int_0^\infty\!\!\int\indic_{\{g(x)>t\}}f(x)\,dx\,dt
=\int f(x)\int_0^{g(x)}dt\,dx=\int g(x)f(x)\,dx.
\]
(The tail formula for $g(X)$ holds for any non-negative random variable, by the
same computation applied to its distribution.) For general $g$ split
$g=g^+-g^-$ and use linearity.
\end{proof}

Linearity of expectation, monotonicity and the variance identities
$\Var(X)=\Exp[X^2]-\mu^2$, $\Var(aX+b)=a^2\Var(X)$ all hold as in the discrete case
(proved in Chapter~\ref{ch:measure} for all random variables simultaneously).

\begin{example}[A distribution with no mean: the Cauchy]\label{ex:cauchy}
The \emph{standard Cauchy} density is $f(x)=\dfrac1{\pi(1+x^2)}$. It integrates to
$1$ (since $\int\frac{dx}{1+x^2}=\arctan x\big|_{-\infty}^\infty=\pi$), but
\[
\int_{-\infty}^\infty\abs x\,f(x)\,dx=\frac2\pi\int_0^\infty\frac x{1+x^2}\,dx
=\frac1\pi\ln(1+x^2)\Big|_0^\infty=\infty .
\]
So $\Exp[X]$ does not exist. The principal-value limit
$\lim_{R\to\infty}\int_{-R}^{R}xf(x)\,dx=0$ exists by symmetry, but this is not an
expectation (it depends on how the two infinite tails are balanced). The Cauchy
law is the standard example of a heavy-tailed distribution; the law of large
numbers fails for it (Chapter~\ref{ch:transforms}).
\end{example}

\section{The standard continuous distributions}

\subsection{Uniform}

$X\sim U(a,b)$ has density $\frac1{b-a}$ on $(a,b)$ and $0$ elsewhere. Then
\[
\Exp[X]=\frac1{b-a}\int_a^bx\,dx=\frac{a+b}2,\qquad
\Exp[X^2]=\frac{b^3-a^3}{3(b-a)}=\frac{a^2+ab+b^2}3,
\]
\[
\Var(X)=\frac{a^2+ab+b^2}3-\frac{(a+b)^2}4=\frac{(b-a)^2}{12}.
\]
Probabilities are proportional to length: $\Prob(X\in(c,d))=(d-c)/(b-a)$.

\subsection{Exponential}

$X\sim\mathrm{Exp}(\lambda)$, $\lambda>0$, has density $\lambda e^{-\lambda x}$ for
$x\ge0$. Then $\Prob(X>x)=e^{-\lambda x}$ and $F(x)=1-e^{-\lambda x}$. By the tail
integral,
\[
\Exp[X]=\int_0^\infty e^{-\lambda t}\,dt=\frac1\lambda,
\]
and integrating by parts, $\Exp[X^2]=\int_0^\infty x^2\lambda e^{-\lambda x}\,dx
=\frac2\lambda\Exp[X]=\frac2{\lambda^2}$, so $\Var(X)=1/\lambda^2$.

\begin{theorem}[Memorylessness of the exponential]\label{thm:expmemoryless}
If $X\sim\mathrm{Exp}(\lambda)$ then $\Prob(X>s+t\given X>s)=\Prob(X>t)$ for all
$s,t\ge0$. Conversely, a positive random variable with a right-continuous
survival function $G(t)=\Prob(X>t)$ that is not identically $0$ or $1$ on
$(0,\infty)$ and satisfies this identity is exponential.
\end{theorem}
\begin{proof}
Forward direction: $e^{-\lambda(s+t)}/e^{-\lambda s}=e^{-\lambda t}$.
Conversely, the identity gives $G(s+t)=G(s)G(t)$. Put $c=G(1)$. By induction
$G(n)=c^n$ and $G(1/m)^m=G(1)$, so $G(n/m)=c^{n/m}$ for all positive rationals. We
have $0<c<1$ (if $c=0$ then $G(1/m)^m=0$ forces $G\equiv0$ on the positive
rationals and, by monotonicity, on $(0,\infty)$; if $c=1$ then $G\equiv1$).
Writing $c=e^{-\lambda}$, $G(r)=e^{-\lambda r}$ for rational $r>0$. Since $G$ is
non-increasing and right-continuous, approximating any real $t$ from the right
by rationals gives $G(t)=e^{-\lambda t}$.
\end{proof}

\begin{proposition}[Exponential and geometric]
If $X\sim\mathrm{Exp}(\lambda)$ then $\lceil X\rceil\sim\mathrm{Geom}(1-e^{-\lambda})$.
\end{proposition}
\begin{proof}
$\Prob(\lceil X\rceil=k)=\Prob(k-1<X\le k)=e^{-\lambda(k-1)}-e^{-\lambda k}=
(e^{-\lambda})^{k-1}(1-e^{-\lambda})$.
\end{proof}

So the exponential is the continuous-time analogue of the geometric distribution:
waiting for a ``memoryless'' event.

\begin{interview}{hazard rate}
The \emph{hazard rate} $h(t)=f(t)/\Prob(X>t)$ is the instantaneous failure rate
given survival to $t$. For the exponential, $h(t)=\lambda$ is constant. A constant
hazard is equivalent to memorylessness; increasing hazard models wear-out, and
decreasing hazard models early failures.
\end{interview}

\subsection{Gamma}

The \emph{gamma function} is $\Gamma(\alpha)=\int_0^\infty x^{\alpha-1}e^{-x}\,dx$
for $\alpha>0$. Integration by parts gives
\[
\Gamma(\alpha+1)=\int_0^\infty x^\alpha e^{-x}\,dx=\bigl[-x^\alpha e^{-x}\bigr]_0^\infty
+\alpha\int_0^\infty x^{\alpha-1}e^{-x}\,dx=\alpha\,\Gamma(\alpha),
\]
and $\Gamma(1)=1$, so $\Gamma(n)=(n-1)!$ for positive integers $n$. Also
$\Gamma(\tfrac12)=\sqrt\pi$ (shown in Section~\ref{sec:normal}).

$X\sim\mathrm{Gamma}(\alpha,\lambda)$ has density
\[
f(x)=\frac{\lambda^\alpha}{\Gamma(\alpha)}x^{\alpha-1}e^{-\lambda x},\qquad x>0 .
\]
It integrates to $1$ by the substitution $u=\lambda x$ in the definition of
$\Gamma(\alpha)$. The same substitution gives the moments:
\[
\Exp[X^k]=\frac{\lambda^\alpha}{\Gamma(\alpha)}\int_0^\infty x^{\alpha+k-1}e^{-\lambda x}\,dx
=\frac{\Gamma(\alpha+k)}{\Gamma(\alpha)\lambda^k}.
\]
Thus $\Exp[X]=\alpha/\lambda$, $\Exp[X^2]=\alpha(\alpha+1)/\lambda^2$ and
$\Var(X)=\alpha/\lambda^2$. For $\alpha=1$ we recover the exponential. For
integer $\alpha=n$ the gamma distribution (the \emph{Erlang} distribution) is the
sum of $n$ independent exponentials, as we prove in
Chapter~\ref{ch:jointcont}.

\subsection{Beta}

$X\sim\mathrm{Beta}(a,b)$, $a,b>0$, has density
\[
f(x)=\frac{x^{a-1}(1-x)^{b-1}}{B(a,b)},\qquad0<x<1,
\]
where $B(a,b)=\int_0^1x^{a-1}(1-x)^{b-1}dx$. In Chapter~\ref{ch:jointcont} we show
$B(a,b)=\Gamma(a)\Gamma(b)/\Gamma(a+b)$, and then
$\Exp[X]=B(a+1,b)/B(a,b)=a/(a+b)$ and $\Var(X)=\dfrac{ab}{(a+b)^2(a+b+1)}$.
The beta family models unknown probabilities, and $\mathrm{Beta}(1,1)=U(0,1)$.

\subsection{The normal distribution}\label{sec:normal}

\begin{theorem}[The Gaussian integral]\label{thm:gauss}
$\displaystyle\int_{-\infty}^\infty e^{-x^2/2}\,dx=\sqrt{2\pi}$.
\end{theorem}
\begin{proof}
Let $I$ be the integral. Since the integrand is non-negative, Tonelli's theorem lets us write
\[
I^2=\int\!\!\int e^{-(x^2+y^2)/2}\,dx\,dy=\int_0^{2\pi}\!\!\int_0^\infty e^{-r^2/2}\,r\,dr\,d\theta
=2\pi\bigl[-e^{-r^2/2}\bigr]_0^\infty=2\pi,
\]
using polar coordinates $(x,y)=(r\cos\theta,r\sin\theta)$ with area element
$r\,dr\,d\theta$. Since $I>0$, $I=\sqrt{2\pi}$.
\end{proof}

Hence $\varphi(x)=\dfrac1{\sqrt{2\pi}}e^{-x^2/2}$ is a density, the \emph{standard
normal} density. We write $Z\sim N(0,1)$, $\Phi(x)=\int_{-\infty}^x\varphi$ for its
CDF (which has no closed form).

\begin{proposition}\label{prop:normalmoments}
If $Z\sim N(0,1)$ then $\Exp[Z]=0$ and $\Var(Z)=1$.
\end{proposition}
\begin{proof}
$x\varphi(x)$ is an odd integrable function, so $\Exp[Z]=0$. By parts, using
$\varphi'(x)=-x\varphi(x)$:
$\Exp[Z^2]=\int x\cdot x\varphi(x)\,dx=\bigl[-x\varphi(x)\bigr]_{-\infty}^\infty
+\int\varphi(x)\,dx=0+1=1$.
\end{proof}

For $\mu\in\R$, $\sigma>0$, the random variable $X=\mu+\sigma Z$ is said to be
$N(\mu,\sigma^2)$. Its density follows from the change of variables rule below:
\[
f_X(x)=\frac1{\sigma\sqrt{2\pi}}\exp\Bigl(-\frac{(x-\mu)^2}{2\sigma^2}\Bigr),
\]
with $\Exp[X]=\mu$ and $\Var(X)=\sigma^2$. Two facts: the family is closed under
affine maps, and standardising $Z=(X-\mu)/\sigma$ reduces every normal
probability to $\Phi$.

\begin{theorem}[Tail bounds for the normal]\label{thm:mills}
For $x>0$,
\[
\Bigl(\frac1x-\frac1{x^3}\Bigr)\varphi(x)\ \le\ \Prob(Z>x)\ \le\ \frac{\varphi(x)}x.
\]
\end{theorem}
\begin{proof}
Upper bound: for $t\ge x$, $1\le t/x$, so
$\Prob(Z>x)=\int_x^\infty\varphi(t)\,dt\le\int_x^\infty\frac tx\varphi(t)\,dt
=\frac{\varphi(x)}x$, since $\frac d{dt}(-\varphi(t))=t\varphi(t)$.
Lower bound: a direct differentiation gives
$\frac d{dt}\bigl[(\tfrac1t-\tfrac1{t^3})\varphi(t)\bigr]=-(1-3t^{-4})\varphi(t)$, hence
$\int_x^\infty(1-3t^{-4})\varphi(t)\,dt=(\tfrac1x-\tfrac1{x^3})\varphi(x)$. As
$(1-3t^{-4})\varphi(t)\le\varphi(t)$, the lower bound follows.
\end{proof}

The bounds show that the normal tail decays like $e^{-x^2/2}$, much faster than
any exponential; this ``thin tail'' is why the normal is so well-behaved and so
dangerous when real data have fat tails.

\section{Functions of a continuous random variable}

\begin{theorem}[Change of variables]\label{thm:changevar}
Let $X$ have density $f_X$, and let $g$ be strictly monotone and differentiable on
an open interval $I$ containing the support of $X$, with inverse $h=g^{-1}$. Then
$Y=g(X)$ has density
\[
f_Y(y)=f_X\bigl(h(y)\bigr)\,\abs{h'(y)}.
\]
\end{theorem}
\begin{proof}
If $g$ is increasing, $F_Y(y)=\Prob(g(X)\le y)=\Prob(X\le h(y))=F_X(h(y))$. Differentiate
by the chain rule: $f_Y(y)=f_X(h(y))h'(y)$ with $h'>0$. If $g$ is decreasing,
$F_Y(y)=\Prob(X\ge h(y))=1-F_X(h(y))$, so $f_Y(y)=-f_X(h(y))h'(y)=f_X(h(y))\abs{h'(y)}$
as $h'<0$.
\end{proof}

When $g$ is not monotone, use the CDF method directly: express $\{g(X)\le y\}$ as
an event about $X$.

\begin{example}[The square of a normal]
Let $Y=Z^2$. For $y>0$, $F_Y(y)=\Prob(-\sqrt y\le Z\le\sqrt y)=2\Phi(\sqrt y)-1$. Differentiating,
\[
f_Y(y)=2\varphi(\sqrt y)\cdot\frac1{2\sqrt y}=\frac1{\sqrt{2\pi}}\,y^{-1/2}e^{-y/2},
\]
which is the Gamma$(\tfrac12,\tfrac12)$ density (the \emph{chi-squared} law with one
degree of freedom). Comparing normalising constants,
$\Gamma(\tfrac12)=\sqrt\pi$.
\end{example}

\begin{theorem}[Probability integral transform]\label{thm:pit}
Let $X$ have a continuous, strictly increasing CDF $F$ on its support. Then
\begin{enumerate}[label=(\roman*)]
\item $F(X)\sim U(0,1)$;
\item if $U\sim U(0,1)$ then $F^{-1}(U)$ has CDF $F$.
\end{enumerate}
\end{theorem}
\begin{proof}
(i) For $u\in(0,1)$, $\Prob(F(X)\le u)=\Prob(X\le F^{-1}(u))=F(F^{-1}(u))=u$, which is
the $U(0,1)$ CDF. (ii) $\Prob(F^{-1}(U)\le x)=\Prob(U\le F(x))=F(x)$.
\end{proof}

Part (ii) is the \emph{inverse transform method} for simulation
(Chapter~\ref{ch:sim}): any distribution can be generated from uniform random
numbers. For example if $U\sim U(0,1)$ then $-\frac1\lambda\ln U\sim\mathrm{Exp}(\lambda)$,
because $F^{-1}(u)=-\frac1\lambda\ln(1-u)$ and $1-U\overset d=U$.

\begin{nuance}{quantiles when the CDF is not strictly increasing}
For general $F$ define the quantile function $F^{-1}(u)=\inf\{x:F(x)\ge u\}$. Then
$F^{-1}(U)$ still has CDF $F$ (because $F^{-1}(u)\le x\iff u\le F(x)$), but
$F(X)$ is uniform only if $F$ is continuous. For discrete $X$, $F(X)$ takes only
finitely or countably many values.
\end{nuance}

\section*{Exercises}
\addcontentsline{toc}{section}{Exercises}

\begin{enumerate}
\item Find $c$ so that $f(x)=cx^2(1-x)$ on $[0,1]$ is a density; compute the mean and
variance.
\item Let $X\sim\mathrm{Exp}(\lambda)$. Find the density of $Y=X^2$ and of
$Y=e^{-X}$ (the latter is uniform).
\item A bus arrives with $\mathrm{Exp}(1/10)$ waiting time (mean 10 minutes). You
have already waited $10$ minutes. What is the expected additional wait?
\item Show that for any random variable with $\Exp[X^2]<\infty$,
$\Exp[X]=\int_0^\infty\Prob(X>t)\,dt-\int_0^\infty\Prob(X<-t)\,dt$.
\item Prove $\Exp[X^n]=\int_0^\infty nt^{n-1}\Prob(X>t)\,dt$ for $X\ge0$.
\item Compute $\Prob(\abs Z>3)$ for $Z\sim N(0,1)$ and compare with the bounds of
Theorem~\ref{thm:mills}. Compute $\Prob(\abs Z>x)$ for the Cauchy.
\item Let $U\sim U(0,1)$. Show that $\min(U,1-U)\sim U(0,\tfrac12)$.
\item Prove that the Cauchy distribution is closed under reciprocals:
if $X$ is standard Cauchy then so is $1/X$.
\item (Challenge) Show that if $X$ has density and $\Exp[X]$ exists then
$\Exp\abs{X-c}$ is minimised at the median, $c=\mathrm{med}(X)$ with
$F(c)=\tfrac12$. (Differentiate $\int\abs{x-c}f(x)\,dx$ in $c$.)
\end{enumerate}
